%! Introduction to Polynomial Functions — Unit 2, Lesson 2.0 — Slides (Beamer)
\documentclass[aspectratio=169, 11pt]{beamer}
\usepackage{precalculus-beamer}   % provides \plumheader{} and \sectionlabel[color]{}

\begin{document}

% TITLE SLIDE
{
\setbeamercolor{background canvas}{bg=plum}
\begin{frame}[plain]
  \vspace{2.8cm}\hspace{0.55cm}%
  \begin{minipage}{0.9\textwidth}
    {\color{goldacc}\bfseries\small LESSON 2.0}\\[0.5cm]
    {\color{white}\bfseries\LARGE Introduction to Polynomial Functions}\\[1.2cm]
    {\color{lilacmid}\small Precalculus~$\cdot$~Unit 2: Polynomial Functions}
  \end{minipage}
\end{frame}
}

% HOOK SLIDE
\begin{frame}[t, plain]
  \plumheader{Hook: Which Three Belong Together?}
  \sectionlabel[goldacc]{Think About It}

  Six expressions. Three belong to the same family; three do not.

  \bigskip
  \begin{center}
    \Large
    $3x^2 - 5$ \qquad $\dfrac{4}{x}$ \qquad $x^3 - 2x + 1$ \qquad
    $\sqrt{x} + 6$ \qquad $7$ \qquad $2^x$
  \end{center}

  \bigskip
  \begin{columns}[T]
    \column{0.48\textwidth}
      \begin{block}{Question 1}
        Which three did you circle --- and what \textbf{rule} did you use?
      \end{block}

    \column{0.48\textwidth}
      \begin{block}{Question 2}
        Does the lone $7$ belong with them, or not? Defend your answer.
      \end{block}
  \end{columns}

  \bigskip
  \textit{Almost everyone sorts these correctly on instinct. Today we sharpen the
  instinct into a definition.}
\end{frame}

% SECTION 1 — WHAT IS A POLYNOMIAL FUNCTION
\begin{frame}[t, plain]
  \plumheader{1. What Is a Polynomial Function?}

  \begin{block}{Informally}
    Add up terms, where every term is a \textbf{number} times $x$ raised to a
    \textbf{whole-number} power.
  \end{block}

  \medskip
  Take $p(x) = 4x^3 - 7x^2 + 2x - 9$ apart: \quad
  $4x^3$, \quad $-7x^2$, \quad $2x$, \quad $-9$.

  \smallskip
  \textit{Quietly, the last two are $2x^1$ and $-9x^0$ --- every term fits the pattern.}

  \bigskip
  \begin{block}{Formally}
    \[
      p(x) = a_n x^n + a_{n-1}x^{n-1} + \cdots + a_1 x + a_0
    \]
    where $n$ is a positive integer, every coefficient is a \textbf{real} number,
    and the leading coefficient $a_n$ is \textbf{not zero}.
  \end{block}

  \medskip
  \textbf{And the lone 7?} A constant is a polynomial too --- of \textbf{degree 0}.
\end{frame}

% SECTION 2 — THE PARTS, BY NAME
\begin{frame}[t, plain]
  \plumheader{2. The Parts of a Polynomial, by Name}

  For $p(x) = 4x^3 - 7x^2 + 2x - 9$:

  \medskip
  \begin{tabular}{|p{0.26\textwidth}|p{0.40\textwidth}|c|}
  \hline
  \textbf{Part} & \textbf{What it is} & \textbf{Here} \\
  \hline
  Degree & the highest exponent on $x$ & ? \\
  \hline
  Leading term & the term carrying that exponent & ? \\
  \hline
  Leading coefficient & the number in front of it & ? \\
  \hline
  Constant term & the term with no $x$ & ? \\
  \hline
  \end{tabular}

  \bigskip
  \begin{block}{Standard form}
    Terms run from the \textbf{highest} power down to the \textbf{lowest} --- which
    parks the leading term and the constant term at the two ends, where they are
    easy to read.
  \end{block}

  \medskip
  \textbf{Watch the sign:} in $-x^3$, the leading coefficient is $\mathbf{-1}$, not
  $1$. The minus belongs to the number.
\end{frame}

% SECTION 3 — SPOTTING THE IMPOSTORS
\begin{frame}[t, plain]
  \plumheader{3. Spotting the Impostors}
  \sectionlabel[redacc]{The Big One}

  Rewrite each one to expose the problem:

  \medskip
  \begin{tabular}{|p{0.22\textwidth}|p{0.24\textwidth}|p{0.40\textwidth}|}
  \hline
  \textbf{Expression} & \textbf{Rewritten} & \textbf{Why it fails} \\
  \hline
  $\dfrac{4}{x}$ & $4x^{-1}$ & the exponent is \textbf{negative} \\
  \hline
  $\sqrt{x} + 6$ & $x^{1/2} + 6$ & the exponent is a \textbf{fraction} \\
  \hline
  $2^x$ & --- & the \textbf{variable} is in the exponent \\
  \hline
  $\dfrac{3}{x+1}$ & --- & the \textbf{variable} is in a denominator \\
  \hline
  \end{tabular}

  \bigskip
  \begin{block}{The three-question test --- all three must be yes}
    \begin{enumerate}\setlength{\itemsep}{3pt}
      \item Is every exponent on $x$ a whole number?
      \item Is $x$ out of every denominator?
      \item Is $x$ out of every exponent?
    \end{enumerate}
  \end{block}
\end{frame}

% SECTION 3, CONT. — THE TWO TRAPS
\begin{frame}[t, plain]
  \plumheader{The Two Traps}

  \begin{columns}[T]
    \column{0.48\textwidth}
      \begin{block}{$\dfrac{x^2}{5}$ \quad IS a polynomial}
        \begin{itemize}\setlength{\itemsep}{6pt}
          \item Dividing by a \textbf{number} is multiplying by $\frac{1}{5}$
          \item So this is just $\frac{1}{5}x^2$
          \item The $x$ is on top, where it belongs
        \end{itemize}
      \end{block}

    \column{0.48\textwidth}
      \begin{block}{$\sqrt{7}\,x^2 - \frac{1}{2}x$ \quad IS a polynomial}
        \begin{itemize}\setlength{\itemsep}{6pt}
          \item The radical is a \textbf{coefficient}, not an exponent
          \item Coefficients may be any real number --- even $\pi$
          \item The exponents here are $2$ and $1$
        \end{itemize}
      \end{block}
  \end{columns}

  \bigskip
  \textbf{In one sentence:} the coefficients are never the problem. Only the
  \textbf{exponents on $x$} --- and where the $x$ is sitting --- decide the question.

  \medskip
  \textit{Challenge:} give me an expression that is one small change away from being
  a polynomial. Now make the change.
\end{frame}

% SECTION 4 — THE FAMILY BY DEGREE
\begin{frame}[t, plain]
  \plumheader{4. The Family, Sorted by Degree}

  \begin{tabular}{|c|c|l|p{0.32\textwidth}|}
  \hline
  \textbf{Degree} & \textbf{Name} & \textbf{Example} & \textbf{Its graph} \\
  \hline
  0 & constant  & $p(x) = 6$              & a flat horizontal line \\
  \hline
  1 & linear    & $p(x) = 2x - 5$         & a straight slanted line \\
  \hline
  2 & quadratic & $p(x) = x^2 - 4x + 3$   & a parabola \\
  \hline
  3 & cubic     & $p(x) = x^3 - 2x$       & new to us this unit \\
  \hline
  4 & quartic   & $p(x) = x^4 - 5x^2 + 4$ & new to us this unit \\
  \hline
  \end{tabular}

  \bigskip
  Degrees 0, 1, and 2 are old friends from Algebra 1, Algebra 2, and Unit 1.

  \medskip
  Lesson 2.1 revisits degree \textbf{2}; the rest of the unit works with degree
  \textbf{3} and higher.
\end{frame}

% SECTION 5 — WHAT THE FAMILY LOOKS LIKE
\begin{frame}[t, plain]
  \plumheader{5. What the Whole Family Looks Like}

  Every polynomial graph is \textbf{smooth} (no corners) and \textbf{unbroken}
  (no holes, no jumps).

  \medskip
  \begin{columns}[T]
    \column{0.52\textwidth}
      \begin{center}
      \begin{tabular}{ccc}
      \begin{tikzpicture}[x=0.42cm, y=0.30cm]
        \draw[->, thick] (-3.2,0) -- (3.2,0);
        \draw[->, thick] (0,-3.2) -- (0,3.4);
        \draw[very thick, plum, domain=-2.8:2.8, samples=80, smooth, variable=\x]
          plot ({\x}, {0.5*\x*\x - 2});
      \end{tikzpicture}
      &
      \begin{tikzpicture}[x=0.42cm, y=0.30cm]
        \draw[->, thick] (-3.2,0) -- (3.2,0);
        \draw[->, thick] (0,-3.2) -- (0,3.4);
        \draw[very thick, plum, domain=-2.6:2.6, samples=80, smooth, variable=\x]
          plot ({\x}, {0.4*\x*\x*\x - 1.6*\x});
      \end{tikzpicture}
      &
      \begin{tikzpicture}[x=0.42cm, y=0.30cm]
        \draw[->, thick] (-3.2,0) -- (3.2,0);
        \draw[->, thick] (0,-3.2) -- (0,3.4);
        \draw[very thick, plum, domain=-2.4:2.4, samples=80, smooth, variable=\x]
          plot ({\x}, {0.35*\x*\x*\x*\x - 1.8*\x*\x + 1});
      \end{tikzpicture}
      \\[2pt]
      \small\textbf{A} (deg 2) & \small\textbf{B} (deg 3) & \small\textbf{C} (deg 4)
      \end{tabular}
      \end{center}

    \column{0.44\textwidth}
      \begin{block}{Turning point}
        A spot where the graph stops rising and starts falling --- or stops falling
        and starts rising.
      \end{block}

      \medskip
      \begin{itemize}\setlength{\itemsep}{5pt}
        \item A: \textbf{1} turn \quad B: \textbf{2} turns \quad C: \textbf{3} turns
        \item Degree $n$ $\Rightarrow$ \textbf{at most} $n - 1$ turning points
        \item Fewer is allowed: $p(x) = x^3$ has none
      \end{itemize}
  \end{columns}

  \bigskip
  \textit{Smooth and unbroken is exactly why polynomials are the family calculus
  starts with.}
\end{frame}

% WHERE THIS UNIT IS GOING
\begin{frame}[t, plain]
  \plumheader{Where This Unit Is Going}

  Today we only \textbf{named} the family. From here:

  \bigskip
  \begin{columns}[T]
    \column{0.48\textwidth}
      \begin{itemize}\setlength{\itemsep}{7pt}
        \item \textbf{2.1} Quadratic Functions Revisited
        \item \textbf{2.2} Polynomial Functions and Rates of Change
        \item \textbf{2.3} Polynomial Functions and Real Zeros
        \item \textbf{2.4} Polynomial Functions and Complex Zeros
      \end{itemize}

    \column{0.48\textwidth}
      \begin{itemize}\setlength{\itemsep}{7pt}
        \item \textbf{2.5} Polynomial Functions and End Behavior
        \item \textbf{2.6} Equivalent Representations of Polynomial Expressions
        \item \textbf{2.7} Polynomial Equations, Inequalities, and Modeling
      \end{itemize}
  \end{columns}

  \bigskip
  Every one of those lessons assumes you can look at an expression and say, without
  hesitating: \textit{``that is a polynomial, and its degree is\ldots''}
\end{frame}

% CLASS PRACTICE
\begin{frame}[t, plain]
  \plumheader{Class Practice}
  \sectionlabel[redacc]{Try It}

  \begin{enumerate}\setlength{\itemsep}{10pt}
    \item For $p(x) = -2x^4 + x^3 - 6$, name the degree, leading term, leading
          coefficient, and constant term. What is this polynomial called, and how
          many turning points could its graph have at most?

    \item Polynomial or not? Give a one-sentence reason for each.

          \medskip
          \begin{columns}[T]
            \column{0.45\textwidth}
              (a) \quad $h(x) = 6x^5 - \dfrac{x}{4} + 11$
            \column{0.45\textwidth}
              (b) \quad $k(x) = 3x^2 + \dfrac{5}{x}$
          \end{columns}

    \item Write $9x - x^3 + 2$ in standard form, then read off its leading
          coefficient. \textit{(Careful with the sign.)}
  \end{enumerate}
\end{frame}

\end{document}
